\section{Transfer Feasibility: What Survives Contact with Animal Signals}
\label{sec:transfer}

We now have two stacks in view: the human translation genealogy
(Section~\ref{sec:humanmt}) and the speech/audio toolkit
(Section~\ref{sec:foundationaudio}). The decisive question is not whether these
tools are powerful---they are---but which of their \emph{presuppositions}
survive when the source is a dog rather than a human. This section audits the
transfer component by component, grounding the audit in the machine-learning
literature on transfer, representation, grounding, and compositionality. The
result, summarized in Table~\ref{tab:transfer}, is a clear pattern:
\emph{the acoustic and representational layers transfer almost intact; the
semantic and evaluative layers do not transfer at all.}

\subsection{Transfer learning: the formal picture}
Transfer learning formalizes when knowledge from a source domain helps in a
target domain~\cite{pan2010transfer,zhuang2021transfer}. Deep features are
often transferable, with generality decreasing as one moves from early to
late layers~\cite{yosinski2014transferable}. This is the optimistic half of
the story: an encoder pretrained on human speech or general audio can be
fine-tuned on barks, and cross-species bioacoustic models already demonstrate
the gain~\cite{hagiwara2023aves,ghani2023perch}. But transfer of
\emph{features} is not transfer of \emph{semantics}. The source domain of human
MT carries a target language with shared reference; the animal target domain
carries no such thing. Feature transfer succeeds precisely where the task is
classification; it cannot manufacture a reference translation that the source
domain never had.

\subsection{Representation-learning foundations}
The machinery that makes transfer possible is the standard deep-learning
stack: backpropagation~\cite{rumelhart1986backprop}, long short-term memory for
sequences~\cite{hochreiter1997lstm}, the deep-learning synthesis of
LeCun, Bengio, and Hinton~\cite{lecun2015deep,bengio2013representation,goodfellow2016deep},
and the optimization and regularization practices (Adam,
dropout) that make large models trainable~\cite{kingma2015adam,srivastava2014dropout}.
These are modality-agnostic: nothing in them assumes human language. A bark
encoder built from the same components is therefore engineerable, and indeed
AVES and NatureLM-audio are built from exactly
them~\cite{hagiwara2023aves,robinson2024naturelm}. The foundations transfer
without remainder; the bottleneck is downstream.

\subsection{The symbol-grounding problem}
Here the transfer breaks. Human translation is grounded: source and target
symbols alike refer to a shared world of objects, actions, and
states~\cite{harnad1990grounding}. A model trained only on text inherits this
grounding indirectly through human usage, but the limits of that inheritance
are well documented---text-only models lack the sensorimotor experience that
gives words their meaning~\cite{bisk2020experience}. For a dog there is no
shared symbolic world to ground into. The dog's signal is grounded in
\emph{its} affective and motivational states, not in a public lexicon; mapping
it to an English sentence is not translation between two grounded languages
but an act of interpretation across an unbridgeable referential gap. This is
the technical content of the ``semantic gap'' thesis, and it is why grounding,
not model scale, is the binding constraint.

\subsection{Compositionality, systematicity, and what probing reveals}
Human language is compositional: the meaning of a whole is a function of the
meanings of its parts and their combination. Whether neural models learn this
is contested~\cite{hupkes2020compositionality}, and diagnostic benchmarks
(SCAN, systematic-generalization tests) show that standard architectures
often fail to generalize compositionally even when they fit the training
distribution~\cite{lake2018scan,lake2023systematic}. Probing studies further
show that what a model ``knows'' is task- and probe-dependent, and that
BERT-style representations encode syntax without necessarily encoding
meaning~\cite{conneau2018cram,rogers2020bertology}. The implication for animals
is twofold. First, even if a bark sequence had compositional structure
(Section~\ref{sec:feasibility}), our models might not recover it. Second, the
absence of compositionality in the \emph{signal}---the likely case for
dogs---means there is no recursive structure for a sequence model to exploit,
so the very property that makes human MT work is missing at the source.

\subsection{Generative and interactive methods: the validation loop}
Two further ML lineages matter not for decoding but for \emph{validation}.
Generative models---GANs~\cite{goodfellow2014gan}, denoising diffusion
models~\cite{ho2020ddpm}, and latent diffusion~\cite{rombach2022ldm}---could
synthesize candidate vocalizations for playback experiments, closing the loop
between hypothesis and animal response. Reinforcement learning---from
DQN~\cite{mnih2015dqn} through PPO~\cite{schulman2017ppo} to the
self-play paradigm of AlphaGo~\cite{silver2016alphago}---supplies the framework
for an interactive agent that learns which inferred intents predict the dog's
behavior. These are the tools that could turn an unfalsifiable interpretation
into a measured, closed-loop system (Section~\ref{sec:eval}). They transfer
well; what they require is an error signal, which the field currently lacks.

We stress that Table~\ref{tab:transfer} is an \emph{analytical framework}, not
an experimental result. Where a row says ``transfers,'' the claim is backed
by measured cross-species fine-tuning gains on BEANS and
BirdSet~\cite{hagiwara2023beans,openbeats2025}; where it says ``fails,'' the
claim rests on the absence of any reference corpus or shared lexicon
(Sections~\ref{sec:eval} and~\ref{sec:lessons}). The framework's value is that
it localizes \emph{where} the bottleneck sits, so effort is not spent scaling
the layers that already transfer.

\begin{table}[htbp]
\centering
\caption{Transfer-feasibility audit: what each layer of the human translation
stack presupposes, and whether that presupposition holds for animal signals.
Verdicts are our assessment.}
\label{tab:transfer}
\small
\setlength{\tabcolsep}{4pt}
\begin{tabular}{@{}p{2.8cm}p{4.4cm}p{4.4cm}p{2.0cm}@{}}
\toprule
\textbf{Component} & \textbf{Human MT presupposes} & \textbf{Animal reality} & \textbf{Verdict} \\
\midrule
Acoustic encoder & raw waveform, any language & raw waveform, any species & transfers \\
Self-supervision & large unlabeled corpus & large unlabeled barks exist & transfers \\
Subword / codec units & discrete reusable units & units discoverable (cf.\ CETI) & partial \\
Sequence model & recursive, compositional source & likely non-compositional signal & limited \\
Parallel corpus & bilingual reference text & no reference exists & fails \\
Grounding & shared public lexicon & private affective states & fails \\
Evaluation metric & ground-truth target sentence & no intended message to score & fails \\
Closed-loop validation & observable error signal & buildable via playback / RL & buildable \\
\bottomrule
\end{tabular}
\end{table}